\begin{table}[htbp]
\centering
\small
\caption{Каналы генерации машинной части: модели, доступ и параметры}
\label{tab:s2}
\begin{tabular}{lllll}
\toprule
Канал & Модель & Оболочка и доступ & Открытые веса & Температура и top-p \\
\midrule
real\_claude & claude-opus-4-8 & вложенный CLI, чистый процесс на ячейку & нет & по умолчанию канала \\
gpt & gpt-5.5 & codex exec, reasoning effort high & нет & по умолчанию канала \\
deepseek\_pro & deepseek-v4-pro & HTTPS API без оболочки & да & 1.0 / 1.0 \\
nemotron & nemotron-3-ultra-free & HTTPS API через шлюз & да & 1.0 / 1.0 \\
\bottomrule
\end{tabular}
\begin{minipage}{\linewidth}\footnotesize\vspace{2pt}
\textit{Примечание.} Единица анализа — канал генерации. Знаменатель: четыре канала, 1079 машинных документов; каждый канал дал 270 документов, кроме nemotron — 269. Данные: реестр корпуса. Таблица описывает условия генерации; чисел прогона в ней нет. Шкала: шкалы нет. Сокращения: top-p — доля вероятностной массы при сэмплировании; «по умолчанию канала» означает, что оболочка параметры не раскрывает. Пропуски: у двух каналов температура и top-p недоступны: оболочка их не сообщает, поэтому межканальные различия нельзя приписывать модели. Поправка на множественные сравнения не применялась.
\end{minipage}
\end{table}
